\section{Simpler examples for the running time}
\label{sec:sparse}

If one only wants inputs on which Algorithm~H does $\Theta(E\sqrt V)$ work,
without the exact maximum number of rounds, a much simpler family will do.
It is the folklore example of disjoint paths with $3,5,7,\ldots$ edges, on
which every phase is forced (compare Hegerfeld and Kratsch's paths
\cite[Section~5.1]{HegerfeldKratsch}), here joined into a tree and given an
input order for Algorithm~H, with exact operation counts.

For $k=2,3,\ldots,K$ let $c_k=k(k-1)/2$ and form a path on the girls
$x_i=g_{c_k+i}$ and boys~$y_i=b_{c_k+i}$, $0\le i<k$:
\[
   x_0\sg y_0\sg x_1\sg y_1\sg\cdots\sg x_{k-1}\sg y_{k-1},
   \qquad\text{with lists}\quad x_0\colon[y_0],\quad x_i\colon[y_{i-1},y_i].
\]
Then, for each $k<K$, append the first boy of component $k+1$ to the end of
the list of the last girl~$x_{k-1}$ of component~$k$. The result is a tree
(in fact a caterpillar) of maximum degree three with
\[
   M=N=s=\frac{K(K+1)}2-1,\qquad E=2s-1 .
\]
For $K=4$, the case of Fig.~\ref{fig:caterpillar}, the lists are
\[
  [1];\ [1,2,3];\ [3];\ [3,4];\ [4,5,6];\ [6];\ [6,7];\ [7,8];\ [8,9].
\]

\begin{figure}[ht]
\centering
\begin{tikzpicture}[scale=1]
\begin{scope}[on background layer]
\draw[sap] (0.000,0.000) -- (0.800,0.000) -- (1.600,0.000) -- (1.600,-0.950);
\end{scope}
\draw[edge,draw=black!55,line width=.7pt] (0.000,0.000) -- (0.800,0.000);
\draw[couple] (1.600,0.000) -- (0.800,0.000);
\draw[edge,draw=black!55,line width=.7pt] (1.600,0.000) -- (1.600,-0.950);
\draw[edge,draw=black!55,line width=.7pt] (1.600,0.000) -- (3.650,0.000);
\draw[edge,draw=black!55,line width=.7pt] (2.850,0.000) -- (3.650,0.000);
\draw[couple] (4.450,0.000) -- (3.650,0.000);
\draw[edge,draw=black!55,line width=.7pt] (4.450,0.000) -- (5.250,0.000);
\draw[couple] (6.050,0.000) -- (5.250,0.000);
\draw[edge,draw=black!55,line width=.7pt] (6.050,0.000) -- (6.050,-0.950);
\draw[edge,draw=black!55,line width=.7pt] (6.050,0.000) -- (8.100,0.000);
\draw[edge,draw=black!55,line width=.7pt] (7.300,0.000) -- (8.100,0.000);
\draw[couple] (8.900,0.000) -- (8.100,0.000);
\draw[edge,draw=black!55,line width=.7pt] (8.900,0.000) -- (9.700,0.000);
\draw[couple] (10.500,0.000) -- (9.700,0.000);
\draw[edge,draw=black!55,line width=.7pt] (10.500,0.000) -- (11.300,0.000);
\draw[couple] (12.100,0.000) -- (11.300,0.000);
\draw[edge,draw=black!55,line width=.7pt] (12.100,0.000) -- (12.100,-0.950);
\node[girl,free] at (0.000,0.000) {1};
\node[girl] at (1.600,0.000) {2};
\node[girl,free] at (2.850,0.000) {3};
\node[girl] at (4.450,0.000) {4};
\node[girl] at (6.050,0.000) {5};
\node[girl,free] at (7.300,0.000) {6};
\node[girl] at (8.900,0.000) {7};
\node[girl] at (10.500,0.000) {8};
\node[girl] at (12.100,0.000) {9};
\node[boy] at (0.800,0.000) {1};
\node[boy,free] at (1.600,-0.950) {2};
\node[boy] at (3.650,0.000) {3};
\node[boy] at (5.250,0.000) {4};
\node[boy,free] at (6.050,-0.950) {5};
\node[boy] at (8.100,0.000) {6};
\node[boy] at (9.700,0.000) {7};
\node[boy] at (11.300,0.000) {8};
\node[boy,free] at (12.100,-0.950) {9};
\end{tikzpicture}

\caption{The tree for $K=4$ after~\Hs1, with the SAP of the first dag round.
Component $k$ is a path whose unique augmenting path has level~$k$; the
thin links between components never enter a dag.}
\label{fig:caterpillar}
\end{figure}

\begin{proposition}
On this tree Algorithm~H performs $K$ rounds. For $j=2,\ldots,K$ the
round at level~$L=j$ finds one SAP, which completes component~$j$, and
\[
   W_5=(2j-1)(K-j+1)+\Iv{j<K},\qquad W_{13}=j,\qquad W_{11}=j(K-j+1).
\]
In total $\sum W_5=K(K-1)(2K+5)/6+K-2$, $\sum W_{13}=K(K+1)/2-1$, and
$\sum W_{11}=K(K-1)(K+4)/6$, while~\Hs1 examines $s$ entries. Hence
$\sum W_5/(E\sqrt V)\to1/3$.
\end{proposition}

\begin{proof}
Step~\Hs1 scans each component's girls~$x_{k-1},\ldots,x_1$ before $x_0$ and
couples $x_i\db y_{i-1}$; it never reaches a link. So each component has a
single free girl~$x_0$ and a single free boy~$y_{k-1}$, its unique augmenting
path has level~$k$, and \A{QUEUE} lists the free girls from component $K$
down to component~2. The round at level~$j$ completes component $j$, whose
free girl is last in the queue, so step~\Hs{16} preserves the order.

The links do not interfere. A link goes from a girl of component $k$ to
the first boy of component $k+1$, so an alternating path, whose steps from
girls to boys follow edges, can only move to components of higher index.
In round~$j$ the free girls lie in components $j,\ldots,K$, so the
completed components $2,\ldots,j-1$ are never entered. The link out of
component $k$ starts at the girl~$x_{k-1}$, which lies at level~$k-1$; in
round~$j$ only levels up to $j-1$ are scanned, so this link is examined
only in round~$k$, and then its boy is already in $B_1$, so step~\Hs6
ignores it. Hence the dags are exactly those of the separate paths.

In round~$j$ each of the $K-j+1$ unresolved components contributes its
girls at levels~$0,\ldots,j-1$: one girl with one neighbour and $j-1$ with
two, $2j-1$ examinations in all, plus one for the link of component~$j$
if $j<K$. The shortest component is scanned last at every level, so the
boys at levels~$1,\ldots,j$ of all unresolved components are marked
before the free boy is found: $W_{11}=j(K-j+1)$. The only root is the free
boy of component~$j$, and its search consumes $j$ arcs.
\end{proof}

Here $\sum W_{13}=\Theta(E)$ but $\sum W_{11}=\Theta(E\sqrt V)$: the mark-resetting loop of
step~\Hs{11} is as expensive as the breadth-first search, and cannot be
left out of a running-time analysis.

\begin{remark}
Replacing component $k$ by a cycle $C_{4k-2}$, with suitable lists, gives a
2-regular graph with $s=K^2-1$ on which every round finds one SAP and
$\sum W_5/(E\sqrt V)\to1/(3\sqrt2)$. The staircase graphs do better on
both counts: they have the largest possible number of rounds, about
$2\sqrt s$, against about $\sqrt{2s}$ for the tree and $\sqrt s$ for the
cycles, and $\sum W_5/(E\sqrt V)\to2\sqrt2/3$
(Theorem~W).
\end{remark}
